We study the fixed-budget best-arm identification (BAI) problem in non-stationary linear bandits.  Concretely, given a fixed time budget $T\in \mathbb{N}$, finite arm set $\mathcal{X} \subset \mathbb{R}^d$, and a potentially adversarial sequence of unknown parameters $\lbrace \theta_t\rbrace_{t=1}^{T}$ (hence non-stationary), a learner aims to identify the arm with the largest cumulative reward $x_* = \arg\max_{x \in \mathcal{X}} x^\top\sum_{t=1}^T \theta_t$ with high probability. In this setting, it is well-known that uniformly sampling arms from the G-optimal design yields a minimax-optimal error probability of $\exp\left(-\Theta\left(T /  H_{G}\right)\right)$, where $H_{G}$ scales proportionally with the dimension $d$. However, this notion of complexity is overly pessimistic, as it is derived from a lower bound in which the arm set consists only of the standard basis vectors, thus masking any potential advantages arising from arm sets with richer geometric structure. To address this, we establish an \textit{arm-set-dependent} lower bound that, in contrast, holds for any arm set. Motivated by the ideas underlying our lower bound, we propose the \textit{Adjacent-optimal design}, a specialization of the well-known $\mathcal{XY}$-optimal design, and develop the \textsf{Adjacent-BAI} algorithm. We prove that the error probability of \textsf{Adjacent-BAI} matches our lower bound up to constants, verifying the tightness of our lower bound, and establishing the arm-set-dependent complexity of this setting.

